% Updated manual appendix draft aligned with the current
% trinity-demonstrator implementation.
%
% This file intentionally does not replace appendix-A.tex.
% It is a separate appendix-style draft that can be included later if wanted.

\chapter{Manual}\label{sec:manual}

\section{Predicates for programming with actors}\label{sec:actor-API}

\noindent \textbf{Predicate:} \texttt{self/1}\hfill\textsc{actor}

\begin{verbatim}
self(-Pid) is det.
\end{verbatim}

\noindent Binds \texttt{Pid} to the process identifier of the calling process.\\

\smallskip
\noindent \textbf{Note:} The runtime uses a global pid model. Conceptually,
pids are of the form \texttt{Id@Node}, although some compatibility-oriented
surfaces, especially certain JSON responses, may still expose local pids as
plain integers.\\

\noindent \textbf{Predicate:} \texttt{spawn/1-3}\hfill\textsc{actor}

\begin{verbatim}
spawn(+Goal) is det.
spawn(+Goal, -Pid) is det.
spawn(+Goal, -Pid, +Options) is det.
\end{verbatim}

\noindent Creates a new Web Prolog actor process running \texttt{Goal}. Valid
options are:

\begin{itemize}
\item
  \texttt{node(+URI)}\\
  Creates the process locally or on a remote compatible node. Default is
  \texttt{localhost}.
\item
  \texttt{monitor(+Boolean)}\\
  If \texttt{true}, monitoring is installed as part of process creation.
  Default is \texttt{false}.
\item
  \texttt{link(+Boolean)}\\
  If \texttt{true}, installs directional parent-to-child link cleanup.
  Default is \texttt{true}.
\item
  \texttt{load\_text(+AtomOrString)}\\
  Loads the clauses specified by a Web Prolog source text into the actor's
  private Prolog database before calling \texttt{Goal}.
\item
  \texttt{load\_list(+ListOfClauses)}\\
  Converts the clauses to source text and loads them into the actor's private
  Prolog database before calling \texttt{Goal}.
\item
  \texttt{load\_uri(+URI)}\\
  Loads the clauses specified by a file path, file URI, or HTTP(S) URI into
  the actor's private Prolog database before calling \texttt{Goal}.
\item
  \texttt{load\_predicates(+ListOfPredicateIndicators)}\\
  Loads the local predicates denoted by
  \texttt{ListOfPredicateIndicators} into the actor's private Prolog database
  before calling \texttt{Goal}.
\end{itemize}

\noindent \textbf{Note:} In \texttt{spawn(+Goal,\,-Pid,\,+Options)}, the option
\texttt{monitor(true)} installs monitoring as part of process creation. When
the child terminates, the parent receives a message
\texttt{down(Ref,\,Pid,\,Reason)} where \texttt{Pid} is the terminated actor,
\texttt{Reason} is its exit reason, and \texttt{Ref} identifies the monitor
instance. For monitoring created via \texttt{monitor(true)}, the current
implementation uses \texttt{Ref\,=\,Pid}.\\

\smallskip
\noindent \textbf{Note:} The semantics of \texttt{link(true)} are directional:
if a parent spawns a child with linking enabled, termination of the parent
causes linked children to be terminated. The link does not imply symmetric
bidirectional exit propagation.\\

\smallskip
\noindent \textbf{Note:} It is possible to pass an arbitrary number of the
\texttt{load\_*} options to \texttt{spawn/3}, and possibly more than one
instance of each variant. To ensure that clauses end up in a well-defined
order, they are converted into Prolog source text before being loaded into the
database. The order of clauses and directives in the resulting source text is
determined by the order of the \texttt{load\_*} options in the option list.\\


\noindent \textbf{Predicate:} \texttt{monitor/2}\hfill\textsc{actor}

\begin{verbatim}
monitor(+PidOrName, -Ref) is det.
\end{verbatim}

\noindent Installs a monitor and returns a fresh reference in \texttt{Ref}. The
first argument may be a pid or a registered local name. When the monitored
process terminates, the monitoring process receives a message of the form
\texttt{down(Ref,\,Pid,\,Reason)}.\\

\noindent \textbf{Predicate:} \texttt{demonitor/1-2}\hfill\textsc{actor}

\begin{verbatim}
demonitor(+Ref) is det.
demonitor(+Ref, +Options) is det.
\end{verbatim}

\noindent Stops monitoring identified by \texttt{Ref}. This is idempotent.
There is only one valid option:

\begin{itemize}
\item
  \texttt{flush}\\
  Non-blocking removal of one pending \texttt{down(Ref,\,\_,\,\_)} message
  from the mailbox.
\end{itemize}

\noindent Using \texttt{monitor(true)} in \texttt{spawn/3} is the safest
variant for short-lived children, since monitoring is established during spawn.
Installing monitoring later with \texttt{monitor/2} is a separate step and may
miss a child that exits immediately.\\

\noindent \textbf{Predicate:} \texttt{register/2}\hfill\textsc{actor}

\begin{verbatim}
register(+Name, +Pid) is det.
\end{verbatim}

\noindent Register a local process under a name, where \texttt{Name} is an atom
and \texttt{Pid} identifies the actor process. The association between the name
and the pid is removed when the process terminates.\\

\noindent \textbf{Predicate:} \texttt{whereis/2}\hfill\textsc{actor}

\begin{verbatim}
whereis(+Name, ?Pid) is det.
\end{verbatim}

\noindent Determine the identity of the actor process associated with the name.
Binds \texttt{Pid} to \texttt{undefined} if the process does not exist.\\

\noindent \textbf{Predicate:} \texttt{unregister/1}\hfill\textsc{actor}

\begin{verbatim}
unregister(+Name) is det.
\end{verbatim}

\noindent Remove the association between the name and the process.\\

\noindent \textbf{Predicate:} \texttt{exit/1}\hfill\textsc{actor}

\begin{verbatim}
exit(+Reason) is det.
\end{verbatim}

\noindent Exit the calling process with \texttt{Reason}.\\

\noindent \textbf{Predicate:} \texttt{exit/2}\hfill\textsc{actor}

\begin{verbatim}
exit(+Pid, +Reason) is det.
\end{verbatim}

\noindent Exit the process identified by \texttt{Pid} with \texttt{Reason}. For
remote actors, the runtime routes this request through the remote node.\\

\noindent \textbf{Predicate:} \texttt{!/2, send/2-3}\hfill\textsc{actor}

\begin{verbatim}
+PidOrName ! +Message is det.
send(+PidOrName, +Message) is det.
send(+PidOrName, +Message, +Options) is det.
\end{verbatim}

\noindent Sends \texttt{Message} to the mailbox of the process identified as
\texttt{PidOrName}. \texttt{PidOrName} may be a local pid, a local registered
name, or a global pid of the form \texttt{Id@Node}. Valid options for
\texttt{send/3} are:

\begin{itemize}
\item
  \texttt{delay(+Number)}\\
  Delays the sending by a specified number of seconds. Default is \texttt{0}.
\item
  \texttt{id(+ID)}\\
  \texttt{ID} is a user-supplied identifier that can be used by
  \texttt{cancel/1} to stop the sending from taking place.
\end{itemize}

\noindent \textbf{Predicate:} \texttt{cancel/1}\hfill\textsc{actor}

\begin{verbatim}
cancel(+ID) is det.
\end{verbatim}

\noindent Tries to cancel the sending of \textit{all} delayed messages with the
specified ID. This is best-effort only, since a message may already have been
sent by the time the call is made.\\

\noindent \textbf{Predicate:} \texttt{output/1-2}\hfill\textsc{actor}

\begin{verbatim}
output(+Data) is det.
output(+Data, +Options) is det.
\end{verbatim}

\noindent Sends a message \texttt{output(Pid,\,Data)} to the target process.
\texttt{Pid} is the pid of the current process. Valid option:

\begin{itemize}
\item
  \texttt{target(+Pid)}\\
  Send the message to \texttt{Pid}. Default is the parent process.
\end{itemize}

\noindent Note that this is just a convenience predicate. A toplevel, like any
other actor, may use \texttt{!/2} to send any term to any process to which it
has a pid.\\

\noindent \textbf{Predicate:} \texttt{input/2-3}\hfill\textsc{actor}

\begin{verbatim}
input(+Prompt, -Data) is det.
input(+Prompt, -Data, +Options) is det.
\end{verbatim}

\noindent Sends a message \texttt{prompt(Pid,\,Prompt)} to the target process and
waits for its input. \texttt{Prompt} may be any non-variable term.
\texttt{Pid} is the pid of the current process. \texttt{Data} will be bound to
the term that the target process sends using \texttt{respond/2}. Valid option:

\begin{itemize}
\item
  \texttt{target(+Pid)}\\
  Send the \texttt{prompt} message to \texttt{Pid}. Default is the parent
  process.
\end{itemize}

\noindent \textbf{Predicate:} \texttt{respond/2}\hfill\textsc{actor}

\begin{verbatim}
respond(+Pid, +Input) is det.
\end{verbatim}

\noindent Sends a response in the form of the term \texttt{Input} to a process
that has prompted its parent process for input.\\

\noindent \textbf{Predicate:} \texttt{receive/1-2}\hfill\textsc{actor}

\begin{verbatim}
receive(+Clauses) is semidet.
receive(+Clauses, +Options) is semidet.
\end{verbatim}

\noindent \texttt{Clauses} is a sequence of receive clauses delimited by a
semicolon:

\begin{verbatim}
{  Pattern1 [if Guard1] ->
         Body1 ;
   ...
   PatternN [if GuardN] ->
         BodyN
}
\end{verbatim}

\noindent Each pattern in turn is matched against the first message in the
mailbox. If a pattern matches and the corresponding guard succeeds, the
matching message is removed from the mailbox and the body of the receive clause
is called. If no message is accepted, the process waits for new messages,
checking them one at a time in arrival order. Messages that are not accepted
are \textit{deferred}, i.e. left in the mailbox without any change in their
contents or order. Valid options:

\begin{itemize}
\item
  \texttt{timeout(+Number)}\\
  If nothing appears in the current mailbox within \texttt{Number} seconds,
  the predicate succeeds anyway. Default is no timeout.
\item
  \texttt{on\_timeout(+Goal)}\\
  If the timeout occurs, \texttt{Goal} is called.
\end{itemize}

\section{Predicates for programming with toplevel actors}\label{sec:toplevel-api}

\noindent \textbf{Predicate:} \texttt{toplevel\_spawn/1-2}\hfill\textsc{actor}

\begin{verbatim}
toplevel_spawn(-Pid) is det.
toplevel_spawn(-Pid, +Options) is det.
\end{verbatim}

\noindent Spawns a toplevel actor and binds \texttt{Pid} to its pid. All options accepted by \texttt{spawn/3} are also accepted by
\texttt{toplevel\_spawn/2}. In addition, \texttt{toplevel\_spawn/2}
accepts the following options:

\begin{itemize}
\item
  \texttt{session(+Boolean)}\\
  If set to \texttt{false}, the toplevel actor terminates after having run a
  goal to completion. If \texttt{true}, further interaction is expected.
  Default is \texttt{false} in the actor API.
\item
  \texttt{target(+Pid)}\\
  Send all answer terms to \texttt{Pid}. Default is the pid of the parent.
\end{itemize}

\noindent \textbf{Predicate:} \texttt{toplevel\_call/2-3}\hfill\textsc{actor}

\begin{verbatim}
toplevel_call(+Pid, +Goal) is det.
toplevel_call(+Pid, +Goal, +Options) is det.
\end{verbatim}

\noindent Asks the toplevel \texttt{Pid} for solutions to \texttt{Goal}. Valid
options are:

\begin{itemize}
\item
  \texttt{template(+Template)}\\
  \texttt{Template} is a term sharing variables with \texttt{Goal}. By
  default, the template is identical to the goal.
\item
  \texttt{offset(+Integer)}\\
  Collect only the slice of solutions starting from \texttt{Integer}. Default
  is \texttt{0}.
\item
  \texttt{limit(+Integer)}\\
  Restrict the length of the returned list of solutions to \texttt{Integer}.
\item
  \texttt{once(+Boolean)}\\
  When \texttt{true}, the list of solutions in a \texttt{success} answer term
  may not be the only solutions that \texttt{Goal} has, but the answer term may
  still indicate that it is. Default is \texttt{false}. Makes sense only in
  combination with the \texttt{limit} option.
\item
  \texttt{target(+Pid)}\\
  Send the answer term to \texttt{Pid}. Default is the value of
  \texttt{target} when passed as an option to \texttt{toplevel\_spawn/2}.
\end{itemize}

\noindent Variables in \texttt{Goal} are not bound directly in the caller.
Instead, solutions and other kinds of output are returned in the form of answer
messages delivered to the mailbox of the target process.

\begin{itemize}
\item
  \texttt{success(Pid, Data, More)}\\
  \texttt{Pid} is the pid of the toplevel process that succeeded in finding
  solutions to \texttt{Goal}. \texttt{Data} is a list holding instantiations of
  \texttt{Template}. \texttt{More} is either \texttt{true} or \texttt{false},
  indicating whether the toplevel can return more solutions if
  \texttt{toplevel\_next/1-2} is called.
\item
  \texttt{failure(Pid)}\\
  \texttt{Pid} is the pid of the toplevel process that failed for lack of
  solutions.
\item
  \texttt{error(Pid, Data)}\\
  \texttt{Pid} is the pid of the toplevel throwing the error.
  \texttt{Data} is the error term.
\end{itemize}


\noindent \textbf{Predicate:} \texttt{toplevel\_next/1-2}\hfill\textsc{actor}

\begin{verbatim}
toplevel_next(+Pid) is det.
toplevel_next(+Pid, +Options) is det.
\end{verbatim}

\noindent Asks toplevel \texttt{Pid} for more solutions. Valid options:

\begin{itemize}
\item
  \texttt{limit(+Integer)}\\
  If omitted, the actor API keeps using the most recent limit value from the
  previous toplevel call state.
\item
  \texttt{target(+Pid)}\\
  Send the answer term to \texttt{Pid}. Default is the value of
  \texttt{target} from the previous call state.
\end{itemize}

\noindent The messages delivered to the target mailbox are the same as for
\texttt{toplevel\_call/2-3}.\\

\noindent \textbf{Predicate:} \texttt{toplevel\_stop/1}\hfill\textsc{actor}

\begin{verbatim}
toplevel_stop(+Pid) is det.
\end{verbatim}

\noindent Asks toplevel \texttt{Pid} to stop searching for more solutions.\\

\noindent \textbf{Predicate:} \texttt{toplevel\_abort/1}\hfill\textsc{actor}

\begin{verbatim}
toplevel_abort(+Pid) is det.
\end{verbatim}

\noindent Tells toplevel \texttt{Pid} to abort the execution of the currently
running goal.\\

\section{Predicates for programming with generic server actors}\label{sec:server-api}

\noindent \textbf{Predicate:} \texttt{server\_spawn/3-4}\hfill\textsc{actor}

\begin{verbatim}
server_spawn(+Pred, +State, -Pid) is det.
server_spawn(+Pred, +State, -Pid, +Options) is det.
\end{verbatim}

\noindent Spawns a generic server actor using callback predicate
\texttt{Pred/4} and initial state \texttt{State}. The callback receives
\texttt{(Request, OldState, Response, NewState)}. All options are forwarded to
\texttt{spawn/3} except:

\begin{itemize}
\item
  \texttt{name(+Name)}\\
  Register the server under \texttt{Name}.
\end{itemize}

\bigskip
\noindent \textbf{Predicate:} \texttt{server\_request/3-4}\hfill\textsc{actor}

\begin{verbatim}
server_request(+To, +Request, -Response) is det.
server_request(+To, +Request, -Response, +Options) is det.
\end{verbatim}

\noindent Makes a synchronous request-response call to the server
\texttt{To}. \texttt{server\_request/4} accepts the options of
\texttt{receive/2}, for example \texttt{timeout(+Seconds)}. Monitoring is
installed automatically so the call fails fast if the server terminates before
replying.\\

\bigskip
\noindent \textbf{Predicate:} \texttt{server\_promise/3-4}\hfill\textsc{actor}

\begin{verbatim}
server_promise(+To, +Request, -Ref) is det.
server_promise(+To, +Request, -Ref, -MonRef) is det.
\end{verbatim}

\noindent Sends \texttt{Request} to the server \texttt{To} and returns a
correlation reference \texttt{Ref}. The reply must later be collected with
\texttt{server\_yield/2-4}. The four-argument variant additionally installs a
monitor and returns its reference in \texttt{MonRef}.\\

\bigskip
\noindent \textbf{Predicate:} \texttt{server\_yield/2-4}\hfill\textsc{actor}

\begin{verbatim}
server_yield(+Ref, -Response) is det.
server_yield(+Ref, -Response, +Options) is det.
server_yield(+Ref, +MonRef, -Response, +Options) is det.
\end{verbatim}

\noindent Waits for the response matching \texttt{Ref}. The three-argument
variant accepts the options of \texttt{receive/2}. The four-argument variant
uses \texttt{MonRef} to detect server termination and throws
\texttt{server\_down(Reason)} if the server dies before replying.\\

\bigskip
\noindent \textbf{Predicate:} \texttt{server\_upgrade/2}\hfill\textsc{actor}

\begin{verbatim}
server_upgrade(+To, +Pred) is det.
\end{verbatim}

\noindent Replaces the callback predicate of the running server \texttt{To}
without disturbing its current state. \texttt{Pred} must be arity 4.\\

\bigskip
\noindent \textbf{Predicate:} \texttt{server\_stop/2}\hfill\textsc{actor}

\begin{verbatim}
server_stop(+To, -Reply) is det.
\end{verbatim}

\noindent Asks the server \texttt{To} to stop gracefully and binds
\texttt{Reply} to its acknowledgement.\\

\section{Predicates for programming with statechart actors}\label{sec:statechart-api}

\noindent \textbf{Predicate:} \texttt{statechart\_spawn/1-2}\hfill\textsc{actor}

\begin{verbatim}
statechart_spawn(-Pid) is det.
statechart_spawn(-Pid, +Options) is det.
\end{verbatim}

\noindent Spawns a statechart actor and binds \texttt{Pid} to its pid.

\medskip
\noindent \texttt{statechart\_spawn/2} requires exactly one source option:

\begin{itemize}
\item
  \texttt{load\_uri(+URI)}\\
  Load and run a statechart from a URI or file path.
\item
  \texttt{load\_text(+Text)}\\
  Load and run a statechart from an XML text string.
\end{itemize}

\noindent All remaining options are passed through to \texttt{spawn/3}.

\medskip
\noindent \texttt{load\_list/1} and \texttt{load\_predicates/1} are rejected
for \texttt{statechart\_spawn/2}.\\

\bigskip
\noindent \textbf{Predicate:} \texttt{raise/1}\hfill\textsc{statechart builtin}

\begin{verbatim}
raise(+Event) is det.
\end{verbatim}

\noindent Enqueues \texttt{Event} on the \textit{internal} event queue of the
current statechart interpreter.

\medskip
\noindent \textbf{Important:} \texttt{raise/1} is only meaningful inside
executable statechart content such as \texttt{<datamodel>},
\texttt{<onentry>}, \texttt{<onexit>}, and \texttt{<go>}.\\

\section{Predicates for programming with supervisor actors}\label{sec:supervisor-api}

\noindent \textbf{Predicate:} \texttt{supervisor\_spawn/2-3}\hfill\textsc{actor}

\begin{verbatim}
supervisor_spawn(+ChildSpecs, -Pid) is det.
supervisor_spawn(+ChildSpecs, -Pid, +Options) is det.
\end{verbatim}

\noindent Spawns a supervisor actor and starts its children. Supported options
are:

\begin{itemize}
\item
  \texttt{strategy(+Strategy)}\\
  One of \texttt{one\_for\_one} (default), \texttt{one\_for\_all}, or
  \texttt{rest\_for\_one}.
\item
  \texttt{intensity(+Integer)}\\
  Max restarts in period. Default \texttt{1}.
\item
  \texttt{period(+Integer)}\\
  Window size for restart intensity. Default \texttt{5}.
\item
  \texttt{name(+Name)}\\
  Register the supervisor under \texttt{Name}.
\end{itemize}

\noindent Other options are passed through to \texttt{spawn/3}.

\medskip
\noindent Child specifications take the form \texttt{child(Id, ChildOptions)}.
Child options are:

\begin{itemize}
\item
  \texttt{start(+Goal)}\\
  Required. Start goal for the child.
\item
  \texttt{restart(+Policy)}\\
  One of \texttt{permanent} (default), \texttt{transient}, or
  \texttt{temporary}.
\item
  \texttt{shutdown(+Shutdown)}\\
  One of \texttt{brutal\_kill}, \texttt{infinity}, or a timeout number.
\item
  \texttt{type(+Type)}\\
  \texttt{worker} (default) or \texttt{supervisor}.
\end{itemize}

\bigskip
\noindent \textbf{Predicate:} \texttt{supervisor\_spawn\_child/3}\hfill\textsc{actor}

\begin{verbatim}
supervisor_spawn_child(+Pid, +ChildSpec, -Reply) is det.
\end{verbatim}

\noindent Dynamically adds and starts one child. \texttt{Reply} is one of
\texttt{ok}, \texttt{error(already\_present)}, or
\texttt{error(start\_failed)}.\\

\bigskip
\noindent \textbf{Predicate:} \texttt{supervisor\_terminate\_child/3}\hfill\textsc{actor}

\begin{verbatim}
supervisor_terminate_child(+Pid, +Id, -Reply) is det.
\end{verbatim}

\noindent Stops child \texttt{Id} but keeps its specification. \texttt{Reply}
is \texttt{ok} or \texttt{error(not\_found)}.\\

\bigskip
\noindent \textbf{Predicate:} \texttt{supervisor\_delete\_child/3}\hfill\textsc{actor}

\begin{verbatim}
supervisor_delete_child(+Pid, +Id, -Reply) is det.
\end{verbatim}

\noindent Removes a non-running child specification. \texttt{Reply} is
\texttt{ok}, \texttt{error(running)}, or \texttt{error(not\_found)}.\\

\bigskip
\noindent \textbf{Predicate:} \texttt{supervisor\_respawn\_child/3}\hfill\textsc{actor}

\begin{verbatim}
supervisor_respawn_child(+Pid, +Id, -Reply) is det.
\end{verbatim}

\noindent Restarts a previously terminated child from its stored specification.
\texttt{Reply} is either \texttt{ok(NewPid)}, \texttt{error(running)},
\texttt{error(not\_found)}, or \texttt{error(start\_failed)}.\\

\bigskip
\noindent \textbf{Predicate:} \texttt{supervisor\_which\_children/2}\hfill\textsc{actor}

\begin{verbatim}
supervisor_which_children(+Pid, -Children) is det.
\end{verbatim}

\noindent Returns a list of:

\begin{verbatim}
info(Id, Pid, Type, Restart)
\end{verbatim}

\noindent where \texttt{Pid} may be \texttt{undefined} for stopped
children.\\

\bigskip
\noindent \textbf{Predicate:} \texttt{supervisor\_count\_children/2}\hfill\textsc{actor}

\begin{verbatim}
supervisor_count_children(+Pid, -Counts) is det.
\end{verbatim}

\noindent Returns:

\begin{verbatim}
[specs-N, active-N, supervisors-N, workers-N]
\end{verbatim}

\bigskip
\noindent \textbf{Predicate:} \texttt{supervisor\_stop/1}\hfill\textsc{actor}

\begin{verbatim}
supervisor_stop(+Pid) is det.
\end{verbatim}

\noindent Stops the supervisor and shuts down children in reverse start
order.\\

\subsubsection{Note on synchronous supervisor calls}

\noindent The dynamic/query API above uses monitored synchronous calls and may
throw:

\begin{itemize}
\item \texttt{supervisor\_down(Reason)}
\item \texttt{supervisor\_call\_timeout(Sup, Request)}
\end{itemize}

\section{Built-in predicates for RPC}\label{sec:synchronous-predicate-api}

\noindent \textbf{Predicate:} \texttt{rpc/2-3}\hfill\textsc{isobase}

\begin{verbatim}
rpc(+URI, +Goal) is nondet.
rpc(+URI, +Goal, +Options) is nondet.
\end{verbatim}

\noindent Semantically, \texttt{rpc/2-3} executes a copy of \texttt{Goal} on
the remote node identified by \texttt{URI} and then unifies the local goal
with the remote copy on backtracking. The current implementation uses the
stateless \texttt{/call} endpoint. The following options are valid:

\begin{itemize}
\item
  \texttt{limit(+Integer)}\\
  Restricts the number of solutions retrieved per roundtrip.
\item
  \texttt{once(+Boolean)}\\
  When \texttt{true}, the returned slice may not be the only solutions that
  \texttt{Goal} has, but the answer term may still indicate that it is.
\item
  \texttt{timeout(+Number)}\\
  Requests a timeout in seconds for goal execution on the remote node. The
  effective timeout is the minimum of the requested value and the timeout
  configured by the node owner.
\item
  \texttt{http\_timeout(+Number)}\\
  Sets the client-side HTTP transport timeout in seconds for each request.
\item
  \texttt{load\_text(+AtomOrString)}\\
  Loads source text into the underlying actor's private database before calling
  \texttt{Goal}.
\item
  \texttt{load\_list(+ListOfClauses)}\\
  Loads a list of clauses into the underlying actor's private database before
  calling \texttt{Goal}.
\item
  \texttt{load\_uri(+URI)}\\
  Loads source text from \texttt{URI} into the underlying actor's private
  database before calling \texttt{Goal}.
\item
  \texttt{load\_predicates(+ListOfPredicateIndicators)}\\
  Loads the local predicates denoted by the supplied predicate indicators into
  the underlying actor's private database before calling \texttt{Goal}.
\end{itemize}

\noindent \texttt{http\_timeout/1} controls only the client-side HTTP transport
timeout. It does not change the remote execution timeout.\\

\noindent \textbf{Predicate:} \texttt{promise/3-4}\hfill\textsc{isobase}

\begin{verbatim}
promise(+URI, +Goal, -Ref) is det.
promise(+URI, +Goal, -Ref, +Options) is det.
\end{verbatim}

\noindent Makes an asynchronous RPC call to node \texttt{URI} with
\texttt{Goal}. The current implementation supports the following options:

\begin{itemize}
\item
  \texttt{template(+Template)}\\
  \texttt{Template} is a term sharing variables with the goal. By default, the
  template is identical to the goal.
\item
  \texttt{offset(+Integer)}\\
  Collect only the slice of solutions starting from \texttt{Integer}. Default
  is \texttt{0}.
\item
  \texttt{limit(+Integer)}\\
  Restrict the number of returned solutions to \texttt{Integer}.
\end{itemize}

\noindent The reference returned in \texttt{Ref} can later be used by
\texttt{yield/2-3} to collect the answer.\\

\noindent \textbf{Predicate:} \texttt{yield/2-3}\hfill\textsc{isobase}

\begin{verbatim}
yield(+Ref, ?Answer) is det.
yield(+Ref, ?Answer, +Options) is det.
\end{verbatim}

\noindent Returns the promised answer from a previous call to
\texttt{promise/3-4}. This predicate must be called by the same process from
which the previous call to \texttt{promise/3-4} was made. Valid options:

\begin{itemize}
\item
  \texttt{timeout(+Number)}\\
  If nothing appears in the current mailbox within \texttt{Number} seconds,
  the goal supplied in the \texttt{on\_timeout} option is called. Default is no
  timeout.
\item
  \texttt{on\_timeout(+Goal)}\\
  If the timeout occurs, \texttt{Goal} is called.
\end{itemize}

\section{Built-in predicates for node deployment}\label{sec:node-predicate-api}

\noindent \textbf{Predicate:} \texttt{node/1-2}\hfill\textsc{isobase}

\begin{verbatim}
node(+Port) is det.
node(+Port, +Options) is det.
\end{verbatim}

\noindent Starts an HTTP server for a Prolog node on \texttt{Port}.

\smallskip
\noindent In the current implementation, the node exposes:

\begin{itemize}
\item \texttt{/}
\item \texttt{/call}
\item \texttt{/toplevel\_spawn}
\item \texttt{/toplevel\_call}
\item \texttt{/toplevel\_next}
\item \texttt{/toplevel\_poll}
\item \texttt{/toplevel\_stop}
\item \texttt{/toplevel\_abort}
\item \texttt{/toplevel\_respond}
\item \texttt{/shell}
\item \texttt{/ws}
\end{itemize}

\noindent The node module defines two owner-controlled settings:

\begin{itemize}
\item
  \texttt{node:cache\_size(+Integer)}\\
  Maximum number of cache entries used by the stateless \texttt{/call}
  endpoint. Default is \texttt{100}.
\item
  \texttt{node:timeout(+Number)}\\
  Owner timeout cap in seconds. Default is \texttt{2}. If a request includes a
  timeout value, the effective timeout is the minimum of that requested value
  and \texttt{node:timeout}.
\end{itemize}

\noindent \texttt{node/2} accepts shared-database startup options:

\begin{itemize}
\item \texttt{load\_shared\_db\_text(+Text)}
\item \texttt{load\_shared\_db\_file(+File)}
\item \texttt{load\_shared\_db\_uri(+URI)}
\end{itemize}

\noindent The shared database is loaded once into a dedicated runtime module
and imported by actor modules.

\section{The stateful WebSocket API}\label{sec:ws-api}

\noindent The WebSocket ACTOR profile lives at \texttt{/ws}. Messages are
encoded as JSON dicts.

\medskip
\noindent Supported commands from client to server are:

\begin{itemize}
\item \texttt{toplevel\_spawn}
\item \texttt{toplevel\_call}
\item \texttt{toplevel\_next}
\item \texttt{toplevel\_stop}
\item \texttt{toplevel\_abort}
\item \texttt{toplevel\_respond}
\item \texttt{spawn}
\item \texttt{send}
\item \texttt{exit}
\end{itemize}

\noindent Supported event types from server to client are:

\begin{itemize}
\item \texttt{spawned}
\item \texttt{success}
\item \texttt{failure}
\item \texttt{error}
\item \texttt{output}
\item \texttt{prompt}
\item \texttt{timeout}
\item \texttt{stop}
\item \texttt{abort}
\item \texttt{responded}
\item \texttt{down}
\end{itemize}

\noindent In WebSocket payloads, pid values may appear as integers for local
actors or as strings of the form \texttt{"Id@Node"} for non-local actors.\\

\section{The semi-stateful HTTP API}\label{sec:semi-stateful-http-api}

\subsubsection{I/O and long-polling}

\noindent The semi-stateful API is built on long-lived toplevel actors plus
per-session queues.

\medskip
\noindent Table~\ref{table:api-parameters-toplevel-spawn} lists the
parameters for the \texttt{/toplevel\_spawn} endpoint.

\begin{table}[h!]
\centering
\begin{tabular}{| m{1.8cm} | m{1.6cm} | m{5.4cm} | m{2.5cm} |}
\hline
\textbf{Parameter} & \textbf{Type} & \textbf{Description} & \textbf{Default Value} \\
\hline
\texttt{format} & atom & Response format, typically \texttt{json} or \texttt{prolog}. & \texttt{json} \\
\hline
\texttt{options} & list or string & Spawn options passed to the underlying toplevel actor. & \texttt{[]} \\
\hline
\texttt{load\_text} & string & Source text to be loaded into the session actor before use. & empty string \\
\hline
\end{tabular}
\caption{HTTP API parameters for the \texttt{/toplevel\_spawn} endpoint.}
\label{table:api-parameters-toplevel-spawn}
\end{table}

\noindent \textbf{Note:} In the actual HTTP ISOTOPE API, \texttt{session(true)}
is always forced by the endpoint, regardless of any \texttt{session/1} value
that may be present in \texttt{options}. The endpoint also injects a small
prelude that redirects \texttt{writeln/1} to \texttt{actor:output/1}.\\

\begin{table}[h!]
\centering
\begin{tabular}{| m{1.6cm} | m{1.4cm} | m{5.6cm} | m{2.7cm} |}
\hline
\textbf{Parameter} & \textbf{Type} & \textbf{Description} & \textbf{Default Value} \\
\hline
\texttt{pid} & pid & The toplevel actor to engage. Accepts integer forms and canonical \texttt{Id@Node} forms. & None \\
\hline
\texttt{goal} & callable & The goal to be called. & None \\
\hline
\texttt{template} & term & The template to be used. For \texttt{format=json}, named query variables are returned as a JSON dict template. & Same as goal \\
\hline
\texttt{offset} & int $>= 0$ & The offset in the virtual list of solutions. & 0 \\
\hline
\texttt{limit} & int $> 0$ & Max number of solutions to be returned. & Effectively no limit \\
\hline
\texttt{format} & atom & Response format. & \texttt{json} \\
\hline
\texttt{load\_text} & string & Additional session-local source text to be loaded before the call. & empty string \\
\hline
\texttt{once} & boolean & Restrict execution to one slice even if more solutions may exist. & \texttt{false} \\
\hline
\texttt{timeout} & number & Requested wait timeout for this call, capped by the node owner timeout. & Owner timeout \\
\hline
\end{tabular}
\caption{HTTP API parameters for the \texttt{/toplevel\_call} endpoint.}
\label{table:api-parameters-toplevel-call}
\end{table}

\begin{table}[h!]
\centering
\begin{tabular}{| m{1.6cm} | m{1.4cm} | m{5.6cm} | m{2.7cm} |}
\hline
\textbf{Parameter} & \textbf{Type} & \textbf{Description} & \textbf{Default Value} \\
\hline
\texttt{pid} & pid & The toplevel actor to continue. Accepts integer forms and canonical \texttt{Id@Node} forms. & None \\
\hline
\texttt{limit} & int $> 0$ & Max number of solutions to be returned by this HTTP wrapper call. & Effectively no limit \\
\hline
\texttt{format} & atom & Response format. & \texttt{json} \\
\hline
\texttt{timeout} & number & Requested wait timeout, capped by the node owner timeout. & Owner timeout \\
\hline
\end{tabular}
\caption{HTTP API parameters for the \texttt{/toplevel\_next} endpoint.}
\label{table:api-parameters-toplevel-next}
\end{table}


\medskip
\noindent \texttt{/toplevel\_poll} takes \texttt{pid}, optional
\texttt{format}, and optional \texttt{timeout}. \texttt{/toplevel\_respond}
takes \texttt{pid}, \texttt{input}, and optional \texttt{format}.
\texttt{/toplevel\_stop} and \texttt{/toplevel\_abort} take \texttt{pid} and
optional \texttt{format}.\\

\medskip
\noindent Table~\ref{table:semi-stateful-node-urls} lists the complete
set of endpoints and their current responses in the semi-stateful API.

\renewcommand{\arraystretch}{1.2}
\begin{table}[h]
\begin{center}
    \begin{tabular}{ | p{4cm} | p{6cm} |}
    \hline
    \textbf{API endpoint} &  \textbf{Response types} \\ \hline
    \texttt{/toplevel\_spawn} & \texttt{spawned} or \texttt{error} \\ \hline
    \texttt{/toplevel\_call} & \texttt{success}, \texttt{failure}, \texttt{error}, \texttt{output}, \texttt{prompt}, \texttt{timeout}, or \texttt{abort} \\ \hline
    \texttt{/toplevel\_next} & \texttt{success}, \texttt{failure}, \texttt{error}, \texttt{output}, \texttt{prompt}, \texttt{timeout}, or \texttt{abort} \\ \hline
    \texttt{/toplevel\_respond} & \texttt{responded} or \texttt{error} \\ \hline
    \texttt{/toplevel\_poll} & \texttt{success}, \texttt{failure}, \texttt{error}, \texttt{output}, \texttt{prompt}, \texttt{timeout}, or \texttt{abort} \\ \hline
    \texttt{/toplevel\_stop} & \texttt{stop} \\ \hline
    \texttt{/toplevel\_abort} & \texttt{abort} \\ \hline
\end{tabular}
\end{center}
\caption{The complete set of endpoints and responses in the semi-stateful API.}
\label{table:semi-stateful-node-urls}
\end{table}

\section{The stateless HTTP API}\label{sec:stateless-http-api}

\noindent The stateless API is exposed at \texttt{/call}.

\begin{table}[h!]
\centering
\begin{tabular}{| m{1.5cm} | m{1.2cm} | m{5.7cm} | m{2.6cm} |}
\hline
\textbf{Parameter} & \textbf{Type} & \textbf{Description} & \textbf{Default Value} \\
\hline
\texttt{goal} & callable & The goal to be called. & None \\
\hline
\texttt{template} & term & The template to be used. For \texttt{format=json}, named query variables are returned as a JSON dict template. & Same as goal \\
\hline
\texttt{offset} & int $>= 0$ & The offset in the virtual list of solutions. & 0 \\
\hline
\texttt{limit} & int $> 0$ & Max number of solutions to be returned. & Effectively no limit \\
\hline
\texttt{load\_text} & string & Source text to be injected into the temporary actor context. & empty string \\
\hline
\texttt{timeout} & number & Requested timeout for remote goal execution, capped by the node owner's timeout setting. & Owner timeout setting \\
\hline
\texttt{once} & boolean-ish atom & Restrict execution to one slice even if more solutions may exist. & \texttt{false} \\
\hline
\texttt{format} & atom & The format of answer responses. & \texttt{json} \\
\hline
\end{tabular}
\caption{HTTP API parameters for the \texttt{/call} endpoint.}
\label{table:api-parameters-call-stateless}
\end{table}

\noindent The only required parameter is \texttt{goal}.

\smallskip
\noindent Although the API is stateless from the client's perspective, the
server may cache a live toplevel actor internally to continue paged solution
retrieval without recomputing earlier work.\\
